\documentclass[10pt, twocolumn]{article}

% ─────────────────────────────────────────────────────────────────
%  CORE PACKAGES
% ─────────────────────────────────────────────────────────────────
\usepackage[T1]{fontenc}
\usepackage[utf8]{inputenc}
\usepackage{lmodern}                  % default CM-style font, clean at small sizes
\usepackage{microtype}                % better justification / kerning
\usepackage[table]{xcolor}
\usepackage{amsmath, amssymb, amsthm}
\usepackage{mathtools}
\usepackage{bm}                       % bold math (\bm{})
\usepackage{graphicx}
\usepackage{float}
\usepackage{booktabs}                 % professional table rules
\usepackage{array}
\usepackage{caption}
\usepackage{subcaption}
\usepackage{hyperref}
\usepackage{cleveref}
\usepackage{enumitem}
\usepackage{algorithm}
\usepackage{algpseudocode}
\usepackage{listings}
\usepackage[skins, breakable, listings]{tcolorbox}
\usepackage{tikz}
\usepackage{eso-pic}
\usetikzlibrary{calc}

% ─────────────────────────────────────────────────────────────────
%  ALTED. COLOR SYSTEM
% ─────────────────────────────────────────────────────────────────
\definecolor{altbg}{HTML}{222222}
\definecolor{altbgraised}{HTML}{2A2A28}
\definecolor{altbgdeep}{HTML}{1A1A18}
\definecolor{alttext}{HTML}{E4E4E4}
\definecolor{altgrey}{HTML}{9E9E9E}
\definecolor{altdg}{HTML}{525252}
\definecolor{altborder}{HTML}{444444}
\definecolor{altborderh}{HTML}{666666}
\definecolor{altdiv}{HTML}{333333}
\definecolor{altshadow}{HTML}{111111}
\definecolor{altlb}{HTML}{888888}
\definecolor{altgreen}{HTML}{6E9A73}
\definecolor{altblue}{HTML}{5F87A1}
\definecolor{altred}{HTML}{A85C5C}
\definecolor{altdkgreen}{HTML}{3A6342}
\definecolor{altdkblue}{HTML}{2D5A7A}
\definecolor{altdkred}{HTML}{8A3535}
\definecolor{altrowave}{HTML}{252523}

% ─────────────────────────────────────────────────────────────────
%  PAGE COLOR & TEXT COLOR
% ─────────────────────────────────────────────────────────────────
\pagecolor{altbg}
\color{alttext}

% ─────────────────────────────────────────────────────────────────
%  BACKGROUND GRID  (very subtle — --div at hairline weight)
% ─────────────────────────────────────────────────────────────────
\AddToShipoutPictureBG{%
  \begin{tikzpicture}[remember picture, overlay,
    help lines/.append style={line width=0.035pt, color=altdiv}]
    \draw[help lines] (current page.south west) grid[step=5pt]
                      (current page.north east);
  \end{tikzpicture}%
}

% ─────────────────────────────────────────────────────────────────
%  GEOMETRY & LAYOUT
% ─────────────────────────────────────────────────────────────────
\usepackage[
    top=1in, bottom=1in,
    left=0.75in, right=0.75in,
    columnsep=0.25in
]{geometry}

% ─────────────────────────────────────────────────────────────────
%  TITLE FORMATTING  (titling + titlesec)
% ─────────────────────────────────────────────────────────────────
\usepackage{titling}
\usepackage{titlesec}
\setlength{\droptitle}{-3em}

\pretitle{%
    \begin{center}
    \large\bfseries\color{alttext}
}
\posttitle{%
    \end{center}
    \vspace{0.2em}
    {\color{altborder}\hrule height 0.5pt}
    \vspace{0.4em}
}
\preauthor{%
    \begin{center}
    \small\color{altgrey}
}
\postauthor{\end{center}}
\predate{%
    \begin{center}
    \small\color{altdg}
}
\postdate{\end{center}\vspace{-0.5em}}

% Section headings — formal, slightly muted, ruled
\titleformat{\section}
  {\normalfont\normalsize\bfseries\color{alttext}}
  {\thesection.}{0.5em}{}[{\color{altborder}\titlerule[0.4pt]}]

\titleformat{\subsection}
  {\normalfont\small\bfseries\color{altgrey}}
  {\thesubsection}{0.5em}{}

\titleformat{\subsubsection}
  {\normalfont\small\itshape\color{altlb}}
  {\thesubsubsection}{0.5em}{}

\titlespacing{\section}{0pt}{1.2ex plus .2ex}{0.8ex plus .1ex}
\titlespacing{\subsection}{0pt}{1ex plus .2ex}{0.5ex plus .1ex}

% ─────────────────────────────────────────────────────────────────
%  CAPTION STYLING
% ─────────────────────────────────────────────────────────────────
\captionsetup{
    font={small, color=altgrey},
    labelfont={bf, color=altlb},
    labelsep=period
}

% ─────────────────────────────────────────────────────────────────
%  HYPERREF COLORS
% ─────────────────────────────────────────────────────────────────
\hypersetup{
    colorlinks=true,
    linkcolor=altblue,
    citecolor=altgreen,
    urlcolor=altblue,
    pdftitle={Implementing the Transformer: Attention Mechanisms and Training Dynamics},
    pdfauthor={Blxke}
}

% ─────────────────────────────────────────────────────────────────
%  LISTINGS  (inline code blocks)
% ─────────────────────────────────────────────────────────────────
\lstdefinestyle{altpython}{
    language=Python,
    basicstyle=\ttfamily\scriptsize\color{alttext},
    keywordstyle=\color{altblue},
    stringstyle=\color{altgreen},
    commentstyle=\color{altdg},
    numberstyle=\tiny\color{altdg},
    numbers=left,
    stepnumber=1,
    numbersep=5pt,
    backgroundcolor=\color{altbgdeep},
    showspaces=false,
    showstringspaces=false,
    showtabs=false,
    tabsize=4,
    breaklines=true,
    frame=none,
    xleftmargin=1em,
    xrightmargin=0.5em,
}

% ─────────────────────────────────────────────────────────────────
%  TCOLORBOX ENVIRONMENTS
% ─────────────────────────────────────────────────────────────────
\AtBeginEnvironment{tcolorbox}{\small}

%% ── Definition box (--blue accent) ────────────────────────────
\newtcolorbox{defn}[1][]{%
    enhanced,
    breakable,
    colback=altbgdeep,
    colframe=altdkblue,
    coltext=alttext,
    coltitle=alttext,
    fonttitle=\bfseries\small,
    title={Definition\ifx&#1&\else\ --- #1\fi},
    arc=0pt, outer arc=0pt,
    leftrule=3pt, rightrule=0.4pt,
    toprule=0.4pt, bottomrule=0.4pt,
    boxsep=4pt, left=4pt, right=4pt,
    top=3pt, bottom=3pt,
    sharp corners
}

%% ── Theorem / Proposition box (--green accent) ────────────────
\newtcolorbox{thm}[1][]{%
    enhanced,
    breakable,
    colback=altbgdeep,
    colframe=altdkgreen,
    coltext=alttext,
    coltitle=alttext,
    fonttitle=\bfseries\small,
    title={Theorem\ifx&#1&\else\ --- #1\fi},
    arc=0pt, outer arc=0pt,
    leftrule=3pt, rightrule=0.4pt,
    toprule=0.4pt, bottomrule=0.4pt,
    boxsep=4pt, left=4pt, right=4pt,
    top=3pt, bottom=3pt,
    sharp corners
}

%% ── Remark / Note (neutral, --border frame) ───────────────────
\newtcolorbox{remark}[1][]{%
    enhanced,
    breakable,
    colback=altbgraised,
    colframe=altborder,
    coltext=alttext,
    coltitle=altgrey,
    fonttitle=\small\itshape,
    title={Remark\ifx&#1&\else: #1\fi},
    arc=0pt, outer arc=0pt,
    boxrule=0.4pt,
    boxsep=4pt, left=4pt, right=4pt,
    top=3pt, bottom=3pt,
    sharp corners
}

%% ── Important / Warning (--red accent) ────────────────────────
\newtcolorbox{important}[1][]{%
    enhanced,
    breakable,
    colback=altbgdeep,
    colframe=altdkred,
    coltext=alttext,
    coltitle=alttext,
    fonttitle=\bfseries\small,
    title={Important\ifx&#1&\else\ --- #1\fi},
    arc=0pt, outer arc=0pt,
    leftrule=3pt, rightrule=0.4pt,
    toprule=0.4pt, bottomrule=0.4pt,
    boxsep=4pt, left=4pt, right=4pt,
    top=3pt, bottom=3pt,
    sharp corners
}

%% ── Code block wrapper ─────────────────────────────────────────
\newtcblisting{codebox}[1][Python]{%
    enhanced,
    breakable,
    colback=altbgdeep,
    colframe=altborder,
    coltext=alttext,
    coltitle=altgrey,
    fonttitle=\ttfamily\scriptsize,
    title={#1},
    arc=0pt, outer arc=0pt,
    boxrule=0.4pt,
    boxsep=3pt, left=2pt, right=2pt,
    top=2pt, bottom=2pt,
    sharp corners,
    listing only,
    listing options={style=altpython}
}

%% ── Equation context box (numbered, titled) ───────────────────
\newtcolorbox{eqbox}[1][]{%
    enhanced,
    colback=altbgraised,
    colframe=altdiv,
    coltext=alttext,
    arc=0pt, outer arc=0pt,
    boxrule=0.4pt,
    boxsep=4pt, left=6pt, right=6pt,
    top=4pt, bottom=4pt,
    sharp corners,
    IfValueT={#1}{\tcbtitle={#1}}
}

% ─────────────────────────────────────────────────────────────────
%  MATH THEOREM ENVIRONMENTS  (unnumbered, use tcolorbox wrappers)
% ─────────────────────────────────────────────────────────────────
\theoremstyle{plain}
\newtheorem{proposition}{Proposition}
\newtheorem{lemma}{Lemma}

% ─────────────────────────────────────────────────────────────────
%  COMMANDS
% ─────────────────────────────────────────────────────────────────
\newcommand{\bb}[1]{\textbf{#1}}
\newcommand{\muted}[1]{\textcolor{altgrey}{#1}}
\newcommand{\acc}[1]{\textcolor{altblue}{#1}}
\newcommand{\code}[1]{\texttt{\textcolor{altgrey}{#1}}}

% Math shorthands
\newcommand{\R}{\mathbb{R}}
\newcommand{\softmax}{\operatorname{softmax}}
\newcommand{\LayerNorm}{\operatorname{LayerNorm}}
\newcommand{\Attn}{\operatorname{Attention}}
\newcommand{\MHA}{\operatorname{MHA}}

% ─────────────────────────────────────────────────────────────────
%  DOCUMENT METADATA
% ─────────────────────────────────────────────────────────────────
\title{%
    CoreML: Implementing the Transformer and the Variations.
}
\author{%
    Bhuvanesh Reddy
}
\date{\small\muted{\today}}

% ════════════════════════════════════════════════════════════════
\begin{document}
% ════════════════════════════════════════════════════════════════

\maketitle
\thispagestyle{empty}

% ─────────────────────────────────────────────────────────────────
%  ABSTRACT
% ─────────────────────────────────────────────────────────────────
\begin{tcolorbox}[
    enhanced,
    colback=altbgraised,
    colframe=altborder,
    coltext=alttext,
    arc=0pt, outer arc=0pt,
    boxrule=0.5pt,
    sharp corners,
    left=6pt, right=6pt, top=5pt, bottom=5pt,
    before upper={\textbf{\small Abstract.}\hspace{0.5em}}
]
\small
I present a from-scratch implementation of the Baseline Transformer architecture along with exploring variants in the positional encoding and attention heads.
\end{tcolorbox}

\vspace{0.5em}

% ─────────────────────────────────────────────────────────────────
\section{Introduction}
% ─────────────────────────────────────────────────────────────────

The Transformer \cite{vaswani2017attention} discards recurrence entirely,
relying solely on attention mechanisms to model dependencies between positions.
This section motivates the architecture and outlines the contributions of this
assignment.

\begin{defn}[Scaled Dot-Product Attention]
Given queries $\bm{Q} \in \R^{n \times d_k}$, keys $\bm{K} \in \R^{m \times d_k}$,
and values $\bm{V} \in \R^{m \times d_v}$, attention is defined as
\begin{equation}
    \Attn(\bm{Q}, \bm{K}, \bm{V}) =
        \softmax\!\left(\frac{\bm{Q}\bm{K}^\top}{\sqrt{d_k}}\right)\bm{V}.
    \label{eq:attention}
\end{equation}
\end{defn}

The $\sqrt{d_k}$ scaling prevents the dot products from entering saturation
regions of the softmax when $d_k$ is large.

% ─────────────────────────────────────────────────────────────────
\section{Architecture}
% ─────────────────────────────────────────────────────────────────

\subsection{Multi-Head Attention}

Rather than computing a single attention function, the model projects $\bm{Q}$,
$\bm{K}$, $\bm{V}$ $h$ times with learned linear projections, attends in
parallel, then concatenates:

\begin{eqbox}
\begin{equation}
    \MHA(\bm{Q},\bm{K},\bm{V}) =
        \operatorname{Concat}(\text{head}_1, \ldots, \text{head}_h)\,\bm{W}^O
    \label{eq:mha}
\end{equation}
\end{eqbox}

where $\text{head}_i = \Attn(\bm{Q}\bm{W}_i^Q,\, \bm{K}\bm{W}_i^K,\,
\bm{V}\bm{W}_i^V)$.

\begin{thm}[Expressivity of MHA]
Multi-head attention with $h$ heads and $d_{\text{model}} = h\,d_k$ parameters
is strictly more expressive than single-head attention with the same parameter
budget, as each head can attend to a distinct subspace of the representation.
\end{thm}

\subsection{Position-wise Feed-Forward Network}

Each encoder/decoder layer includes a two-layer FFN applied identically to
each position:
\begin{equation}
    \operatorname{FFN}(\bm{x}) = \max(0,\, \bm{x}\bm{W}_1 + \bm{b}_1)\bm{W}_2 + \bm{b}_2
\end{equation}

\begin{remark}[Dimensionality]
The inner dimension is typically $d_{ff} = 4\,d_{\text{model}}$.
In the base model \cite{vaswani2017attention}, $d_{\text{model}} = 512$
and $d_{ff} = 2048$.
\end{remark}

\subsection{Positional Encoding}

Since attention is permutation-equivariant, position information must be
injected explicitly. Sinusoidal encodings are given by:
\begin{align}
    \text{PE}_{(pos,\,2i)}   &= \sin\!\left(pos / 10000^{2i/d}\right) \\
    \text{PE}_{(pos,\,2i+1)} &= \cos\!\left(pos / 10000^{2i/d}\right)
\end{align}

\begin{important}[Learned vs.\ Sinusoidal]
Learned positional embeddings and sinusoidal encodings achieve comparable
performance; however, sinusoidal encodings generalise to sequence lengths
unseen during training.
\end{important}

% ─────────────────────────────────────────────────────────────────
\section{Implementation}
% ─────────────────────────────────────────────────────────────────

\subsection{Attention Module}

The core attention function in PyTorch:

\begin{codebox}[Python — Scaled Dot-Product Attention]
import torch
import torch.nn.functional as F
import math

def scaled_dot_product_attention(Q, K, V, mask=None):
    """
    Q, K, V : (batch, heads, seq, d_k)
    """
    d_k = Q.size(-1)
    scores = torch.matmul(Q, K.transpose(-2, -1)) / math.sqrt(d_k)
    if mask is not None:
        scores = scores.masked_fill(mask == 0, float('-inf'))
    weights = F.softmax(scores, dim=-1)
    return torch.matmul(weights, V), weights
\end{codebox}

\subsection{Layer Normalisation}

Residual connections are followed by layer normalisation:
\begin{equation}
    \bm{y} = \LayerNorm\!\left(\bm{x} + \operatorname{Sublayer}(\bm{x})\right)
\end{equation}

\begin{remark}[Pre-LN vs.\ Post-LN]
Modern implementations (GPT, LLaMA) typically apply normalisation
\emph{before} the sublayer (\emph{Pre-LN}), which improves gradient flow and
stabilises training without a learning rate warm-up.
\end{remark}

% ─────────────────────────────────────────────────────────────────
\section{Training Dynamics}
% ─────────────────────────────────────────────────────────────────

\subsection{Learning Rate Schedule}

The original Transformer uses a warm-up schedule:
\begin{equation}
    \alpha = d_{\text{model}}^{-0.5}
              \cdot \min\!\left(\mathit{step}^{-0.5},\;
              \mathit{step} \cdot \mathit{warmup}^{-1.5}\right)
\end{equation}

\begin{important}[Warm-up Critical]
Without the linear warm-up phase the model often diverges in early training.
A typical value is $\mathit{warmup} = 4000$ steps.
\end{important}

\subsection{Regularisation}

Three regularisation strategies are employed:

\begin{enumerate}[leftmargin=*, itemsep=2pt, topsep=3pt]
    \item \textbf{Dropout} applied to attention weights and sublayer outputs
          ($p_{\text{drop}} = 0.1$).
    \item \textbf{Label smoothing} with $\varepsilon_{ls} = 0.1$, which hurts
          perplexity but improves BLEU.
    \item \textbf{Weight decay} ($\lambda = 10^{-4}$) on all non-bias parameters.
\end{enumerate}

% ─────────────────────────────────────────────────────────────────
\section{Experimental Results}
% ─────────────────────────────────────────────────────────────────

\subsection{Ablation: Number of Heads}

\begin{table}[H]
\centering
\small
\setlength{\tabcolsep}{6pt}
\begin{tabular}{lccc}
\toprule
\rowcolor{altbgraised}
\textbf{Heads} ($h$) & \textbf{Val.\ Loss} & \textbf{BLEU} & \textbf{Params} \\
\midrule
1  & 3.42 & 18.1 & 37.4M \\
\rowcolor{altrowave}
4  & 2.91 & 23.7 & 37.4M \\
8  & 2.78 & 25.3 & 37.4M \\
\rowcolor{altrowave}
16 & 2.81 & 24.9 & 37.4M \\
\bottomrule
\end{tabular}
\caption{Effect of the number of attention heads (base model, $d_{\text{model}}=512$).}
\label{tab:heads}
\end{table}

Eight heads yields the best BLEU score; increasing to 16 slightly degrades
performance, consistent with findings in \cite{voita2019analyzing}.

% ─────────────────────────────────────────────────────────────────
\section{Conclusion}
% ─────────────────────────────────────────────────────────────────

We have implemented and analysed the Transformer architecture from scratch,
covering attention, positional encodings, layer normalisation, and training
dynamics. Ablation results confirm the importance of multi-head attention and
proper learning rate scheduling. Future work will explore relative positional
encodings and sparse attention variants.

% ─────────────────────────────────────────────────────────────────
%  REFERENCES  (replace with .bib file as needed)
% ─────────────────────────────────────────────────────────────────
\begin{thebibliography}{9}
\small
\bibitem{vaswani2017attention}
A.\ Vaswani et al., ``Attention is all you need,''
\textit{NeurIPS}, 2017.

\bibitem{voita2019analyzing}
E.\ Voita et al., ``Analyzing multi-head self-attention: Specialized heads
do the heavy lifting, the rest can be pruned,''
\textit{ACL}, 2019.
\end{thebibliography}

\end{document}